\chapter{Desain Sistem}
\label{bab:desain-sistem}

Uraian pada bab ini meliputi model yang digunakan, rancangan tesis, variabel
dalam tesis, teknik pengumpulan data dan analisis data. Awali pembahasan pada
bab ini dengan penjelasan umum tentang solusi yang ditawarkan untuk menjawab
Problem. Akan lebih mudah dicerna, apabila penjelasan ini disertai dengan
diagram sistem secara high-level view sehingga pembaca mendapatkan gambaran
menyeluruh tentang desain sistem untuk menyelesaikan Problem. Setelah itu,
penulis dapat menguraikan desain sistem yang digunakan dalam tesis secara
rinci. Berikut contoh diagram desain sistem secara high-level view dan contoh
sistematika pembahasan dari diagram desain.

\begin{figure}[h]
  \centering
  % Ganti dengan \includegraphics[width=...]{figures/nama-file}
  \begin{tikzpicture}
    \draw[thick] (0,0) rectangle (89.9mm,71.3mm);
    \foreach \i/\t in {0/Bagian 1,1/Bagian 2,2/Bagian 3} {
      \node[draw,thick,minimum width=60mm,minimum height=14mm]
        at (44.95mm,58mm-\i*22mm) {\t};}
    \draw[-{Latex},thick] (44.95mm,51mm) -- (44.95mm,43mm);
    \draw[-{Latex},thick] (44.95mm,29mm) -- (44.95mm,21mm);
  \end{tikzpicture}
  \caption{Desain sistem dari solusi yang ditawarkan}
  \label{fig:desain-sistem}
\end{figure}

Gambar diletakkan rata tengah, dengan menyisakan 1 (satu) baris kosong diatas
dan dibawah gambar. Setiap gambar harus mempunyai nomer identitas gambar dan
diiringi dengan keterangan gambar, yang dituliskan rata tengah dan tebal. Nomer
identitas dan keterangan gambar dituliskan pada 1 (satu) baris dibawah gambar.
Nomer identitas terdiri dari nomer bab dan nomer urutan gambar pada bab
tersebut. Setiap gambar harus dirujuk dan dibahas pada pembahasan dalam
paragraf, seperti kalimat berikut. Gambar~\ref{fig:desain-sistem} menunjukkan
bagan desain sistem yang mempunyai tiga bagian. Jika isi gambar adalah kutipan,
maka penulis dapat menyebutkan sumber referensi dari gambar dibawah gambar dan
diatas identitas gambar, dengan rata tengah dan ditulis dengan ukuran 10 point,
seperti yang terlihat pada Gambar~\ref{fig:kutipan}.

\begin{figure}[h]
  \centering
  \fbox{\parbox[c][58.5mm][c]{78.6mm}{\centering Gambar hasil kutipan}}
  \sumber{http://cdn3.techworld.com/cmsdata/features/3210134/cisco-search-patent.jpg}
  \caption{Contoh gambar kutipan}
  \label{fig:kutipan}
\end{figure}

\section{Bagian 1}

Disini penulis dapat menjelaskan lebih terperinci apa saja yang ada pada bagian
ini. Jika bagian ini mempunyai sub bagian yang perlu diperjelas dalam
pembahasan, penulis dapat menuliskannya dalam sub pembahasan pada bagian ini.

\subsection{Aaaa}

Disini penulis dapat membahas sub bagian Aaaa lebih terperinci. Deskripsi
pembahasan seharusnya singkat, padat dan jelas, sehingga membuat pembaca
memahami maksud penulis yang tertuang dalam tulisan. Apabila pembahasan penulis
memerlukan penulisan persamaan matematis, penulis dapat menuliskannya seperti
dibawah ini.
\begin{equation}
  y_t = y_{t-1} + \alpha \cdot \left(x_t - y_{t-1}\right)
  \label{eq:contoh}
\end{equation}
Penulisan persamaan diletakkan pada baris sendiri rata kiri yang masuk 10 mm
dari batas kiri, dengan menyisakan 1 (satu) baris kosong diatas dan dibawah
gambar. Setiap persamaan harus mempunyai nomer identitas persamaan yang
dituliskan rata kanan dan tebal. Setiap persamaan harus dirujuk dan dibahas
pada pembahasan dalam paragraf, seperti kalimat berikut.
Persamaan~\eqref{eq:contoh} menunjukkan keterhubungan antara fungsi pada waktu
sekarang dan sebelumnya.

Untuk cara penulisan tabel, tabel diletakkan rata tengah, dengan menyisakan 1
(satu) baris kosong diatas dan dibawah tabel. Setiap tabel harus mempunyai
nomer identitas tabel dan diiringi dengan keterangan tabel, yang dituliskan
rata tengah dan tebal. Nomer identitas dan keterangan tabel dituliskan pada 1
(satu) baris dibawah tabel. Nomer identitas terdiri dari nomer bab dan nomer
urutan tabel pada bab tersebut. Setiap tabel harus dirujuk dan dibahas pada
pembahasan dalam paragraf, seperti kalimat berikut. Tabel~\ref{tab:contoh}
menunjukkan contoh penulisan tabel, yang terdiri dari nomer identitas dan
keterangan tabel, dan kemudian isi tabel. Judul pada tabel dapat dituliskan
rata tengah, tebal dan berlatar-belakang agak gelap. Jika isi tabel adalah
kutipan, maka penulis dapat menyebutkan sumber referensi dari tabel dibawah
tabel dengan rata tengah dan ditulis dengan ukuran 10 point, seperti yang
terlihat pada Tabel~\ref{tab:contoh}. Satu tabel tidak boleh melebihi dari 1
(satu) halaman. Jika isi tabel terlalu banyak lebih dari 1 (satu) halaman,
penulis dapat memecah tabel dan memberikan identitas tabel yang berbeda.

\begin{table}[h]
  \caption{Contoh penulisan tabel}
  \label{tab:contoh}
  \centering
  \begin{tabular}{|C{45.9mm}|C{18.3mm}|C{23.6mm}|}
    \hline
    \headrow \textbf{Kolom 1} & \textbf{Kolom 2} & \textbf{Kolom 3} \\
    \hline
     & & \\ \hline
     & & \\ \hline
     & & \\ \hline
     & & \\ \hline
     & & \\ \hline
  \end{tabular}
  \sumber{Badan Pusat Pengolahan Data, 2012 \cite{pengarang2012seminar}}
\end{table}

\subsection{Bbbb}

Disini penulis dapat membahas sub bagian Bbbb lebih terperinci. Deskripsi
pembahasan seharusnya singkat, padat dan jelas, sehingga membuat pembaca
memahami maksud penulis yang tertuang dalam tulisan.

\section{Bagian 2}

Disini penulis dapat menjelaskan lebih terperinci apa saja yang ada pada
Bagian 2 ini. Jika bagian ini mempunyai sub bagian yang perlu diperjelas dalam
pembahasan, penulis dapat menuliskannya dalam sub pembahasan pada bagian ini.

\subsection{Cccc}

Disini penulis dapat membahas sub bagian Cccc lebih terperinci. Deskripsi
pembahasan seharusnya singkat, padat dan jelas, sehingga membuat pembaca
memahami maksud penulis yang tertuang dalam tulisan.

\subsection{Dddd}

Disini penulis dapat membahas sub bagian Dddd lebih terperinci. Deskripsi
pembahasan seharusnya singkat, padat dan jelas, sehingga membuat pembaca
memahami maksud penulis yang tertuang dalam tulisan.

\section{Bagian 3}

Disini penulis dapat menjelaskan lebih terperinci apa saja yang ada pada
Bagian 3 ini. Jika bagian ini mempunyai sub bagian yang perlu diperjelas dalam
pembahasan, penulis dapat menuliskannya dalam sub pembahasan pada bagian ini.

\subsection{Eeee}

Disini penulis dapat membahas sub bagian Eeee lebih terperinci. Deskripsi
pembahasan seharusnya singkat, padat dan jelas, sehingga membuat pembaca
memahami maksud penulis yang tertuang dalam tulisan.

\subsection{Ffff}

Disini penulis dapat membahas sub bagian Ffff lebih terperinci. Deskripsi
pembahasan seharusnya singkat, padat dan jelas, sehingga membuat pembaca
memahami maksud penulis yang tertuang dalam tulisan.
